\subsection{深度的界}

命题 12.--- 设 A 为 Noether 局部环，M 为非零的有限生成 A-模，J 为异于 A 的 A 的理想。我们有不等式链

\[
\begin{array}{r l} \operatorname{prof} _ {A} (J; M) & \leqslant \operatorname{codim} (\operatorname{Supp} (M) \cap V (J), \operatorname{Supp} (M)) \leqslant \dim (M) - \dim (M / J M) \\ & \leqslant [ J / \mathfrak {m} _ {A} J: \kappa_ {A} ]. \end{array}
\]

对每个 \(\mathfrak{p} \in \operatorname{Supp}(\mathrm{M})\cap \mathrm{V}(\mathrm{J})\) ，\(\operatorname{prof}_{\mathrm{A}}(\mathrm{J};\mathrm{M})\) 小于 \(\dim_{\mathrm{A}_{\mathfrak{p}}}(\mathrm{M}_{\mathfrak{p}})\) （第 5 节, 命题 8 及第 4 节, 定理 2 的推论 2），即 (VIII, § 1, 第 4 节, 命题 9)

即 \(\operatorname{codim}(V(\mathfrak{p}),\operatorname{Supp}(\mathrm{M}))\) 。当 \(\mathfrak{p}\) 遍历 \(\operatorname{Supp}(\mathrm{M}) \cap \mathrm{V}(\mathrm{J})\) 时，\(V(\mathfrak{p})\) 遍历 \(\operatorname{Supp}(\mathrm{M}) \cap \mathrm{V}(\mathrm{J})\) 的不可约闭子集，由此得第一个不等式。第二个不等式由 VIII, § 1, 第 2 节的命题 3 得出。此外，可以找到 J 的一个基数为 \([J/\mathfrak{m}_{\mathrm{A}}J:\kappa_{\mathrm{A}}]\) 的生成组 (II, § 3, 第 2 节, 命题 4 的推论 2)；第三个不等式于是由 VIII, § 3, 第 2 节, 公式 (8) 得出。

注 1.-- 考察命题 12 中的不等式链。

\begin{enumerate}
\def\labelenumi{\alph{enumi})}
\item
  要使 \(\mathrm{prof}_{\mathrm{A}}(\mathrm{J};\mathrm{M}) = [\mathrm{J} / \mathfrak{m}_{\mathrm{A}}\mathrm{J}:\kappa_{\mathrm{A}}]\) 成立，必须且只需 J 可由一个 M-正则序列生成（第 3 节, 定理 1 的推论 2 及 A, X, 第 160 页, 推论 1）。
\item
  等式 \(\dim(\mathrm{M}) - \dim(\mathrm{M}/\mathrm{JM}) = [\mathrm{J}/\mathfrak{m}_{\Lambda}\mathrm{J} : \kappa_{\Lambda}]\) 意味着 J 可由一个对 M 正则的序列生成 (VIII, § 3, 第 2 节)。
\item
  若 M 是 Macaulay 的，则有 \(\operatorname{prof}_{\mathrm{A}}(\mathbf{J};\mathbf{M}) = \dim (\mathbf{M}) - \dim (\mathbf{M} / \mathbf{JM})\) (§ 2, 第 2 节, 命题 2 的推论)。
\end{enumerate}

引理 3.--- 设 A 为 Noether 环，\(\mathfrak{p} \subset \mathfrak{p}_1 \subset \ldots \subset \mathfrak{p}_{r-1} \subset \mathfrak{q}\) 为 A 的素理想组成的长度为 \(r\) 的饱和链，M 为有限型 \(\Lambda\) -模，\(n\) 为整数。若 A-模 \(\mathrm{Ext}_{\mathrm{A}_{\mathfrak{p}}}^{n}(\kappa(\mathfrak{p}), \mathrm{M}_{\mathfrak{p}})\) 非零，则 \(\mathrm{Ext}_{\mathrm{A}_{\mathfrak{q}}}^{n+r}(\kappa(\mathfrak{q}), \mathrm{M}_{\mathfrak{q}})\) 亦然。

显然只需处理 \(r = 1\) 的情形；然后把 A、M、p 与 q 分别用 \(A_q\) 、 \(M_q\) 、 \(pA_q\) 与 \(qA_q\) 代换，便归结为处理 A 为局部环且 \(q = m_A\) 的情形。设 \(x\) 为 \(m_A - p\) 的元素。\(A_p\) -模 \(\text{Ext}_A^n(A/p, M) \otimes_A A_p\) 同构于 \(\text{Ext}_{A_p}^n(\kappa(p), M_p)\) (A, X, 第 111 页, 命题 10 b))，故按假设非零；更有 \(\text{Ext}_A^n(A/p, M)\) 不为零。正合序列

\[
0 \rightarrow \mathrm{A} / \mathfrak {p} \stackrel {{x _ {\mathrm{A} / \mathfrak {p}}}} {{\longrightarrow}} \mathrm{A} / \mathfrak {p} \longrightarrow \mathrm{A} / (\mathfrak {p} + x \mathrm{A}) \rightarrow 0
\]

给出扩张模的一个正合序列

\[
\operatorname{Ext} _ {\mathrm{A}} ^ {n} (\mathrm{A} / \mathfrak {p}, \mathrm{M}) \xrightarrow {u} \operatorname{Ext} _ {\mathrm{A}} ^ {n} (\mathrm{A} / \mathfrak {p}, \mathrm{M}) \longrightarrow \operatorname{Ext} _ {\mathrm{A}} ^ {n + 1} (\mathrm{A} / (\mathfrak {p} + x \mathrm{A}), \mathrm{M}),
\]

其中 \(u\) 是乘以 \(x\) 的映射。由 Nakayama 引理，它不是满射，故 A-模 \(\text{Ext}_A^{n+1}(\text{A}/(\mathfrak{p} + x\text{A}), \text{M})\) 不为零。而 A 的包含 \(\mathfrak{p} + x\text{A}\) 的素理想只有 \(\mathfrak{m}_\text{A}\) ；因此 A-模 \(\text{A}/(\mathfrak{p} + x\text{A})\) 有有限长度 (VIII, § 3, 第 2 节, 引理 2)。若 \(\text{Ext}_\text{A}^{n+1}(\kappa_\text{A}, \text{M})\) 为零，则由对 N 的长度作归纳将推出，对每个有限长度的 A-模 N，\(\text{Ext}_\text{A}^{n+1}(\text{N}, \text{M})\) 为零。这一矛盾完成了证明。

命题 13.--- 设 A 为 Noether 环，M 为有限生成 A-模，p 与 q 为 Supp(M) 中满足 p ⊂ q 的两个素理想。则

\[
\operatorname{prof} _ {A _ {\mathfrak {q}}} (M _ {\mathfrak {q}}) \leqslant \operatorname{prof} _ {A _ {\mathfrak {p}}} (M _ {\mathfrak {p}}) + \dim (A _ {\mathfrak {q}} / \mathfrak {p} A _ {\mathfrak {q}}).
\]

更精确地说，对素理想的每个饱和链 \(\mathfrak{p} \subset \mathfrak{p}_1 \subset \ldots \subset \mathfrak{p}_{r-1} \subset \mathfrak{q}\) ，有 \(\operatorname{prof}_{\mathrm{A}_{\mathfrak{q}}}(\mathrm{M}_{\mathfrak{q}}) \leqslant \operatorname{prof}_{\mathrm{A}_{\mathfrak{p}}}(\mathrm{M}_{\mathfrak{p}}) + r\) 。

令 \(p = \text{prof}_{\text{A}_p}(\text{M}_p)\) ，我们来证明第二个不等式。若 \(p = +\infty\) 它是显然的；否则，有 \(\text{Ext}_{\text{A}_p}^p(\kappa(\mathfrak{p}), \text{M}_p) \neq 0\) ，从而由引理 3 得 \(\operatorname{Ext}_{\Lambda_{\mathfrak{q}}}^{p+r}(\kappa(\mathfrak{q}),\mathrm{M}_{\mathfrak{q}})\neq0\) ，这蕴含 \(\operatorname{prof}_{\mathrm{A}_{\mathfrak{q}}}(M_{\mathfrak{q}})\leqslant p+r\) 。由于 \(\dim(A_{\mathfrak{q}}/\mathfrak{p}A_{\mathfrak{q}})\) 是以给定端点 \(\mathfrak{p}\) 与 \(\mathfrak{q}\) 的素理想饱和链的长度的上确界，第一个断言是第二个断言的推论。

推论 1.--- 我们有不等式

\[
\dim (M _ {\mathfrak {q}}) - \operatorname{prof} _ {\Lambda_ {\mathfrak {q}}} (M _ {\mathfrak {q}}) \geqslant \dim (M _ {\mathfrak {p}}) - \operatorname{prof} _ {A _ {\mathfrak {p}}} (M _ {\mathfrak {p}}) \geqslant 0.
\]

这由命题 13 及不等式 \(\dim(M_{\mathfrak{q}}) \geqslant \dim(M_{\mathfrak{p}}) + \dim(A_{\mathfrak{q}} / \mathfrak{p}A_{\mathfrak{q}})\) (VIII, § 1, 第 4 节, 命题 9, a) 及第 3 节, 命题 7, b)) 得出。

推论 2.--- 设 A 为 Noether 局部环，M 为有限生成 A-模。则有不等式

\[
\operatorname{prof} _ {\mathrm{A}} (\mathrm{M}) \leqslant \inf _ {\mathfrak {p} \in \operatorname{Ass} (\mathrm{M})} \dim (\mathrm{A} / \mathfrak {p}).
\]

设 \(\mathfrak{p}\) 为与 M 相伴的素理想；有 \(\mathrm{prof}_{\mathrm{A}_{\mathfrak{p}}}(\mathrm{M}_{\mathfrak{p}})=0\) （第 1 节, 注 2）。把命题 13 应用于理想 \(\mathfrak{p}\subset\mathfrak{m}_{\mathrm{A}}\) 便给出不等式 \(\mathrm{prof}_{\mathrm{A}}(\mathrm{M})\leqslant\dim(\mathrm{A}/\mathfrak{p})\) ，从而得本推论。

注 2.--- 我们有 \(\sup_{\mathfrak{p}\in\operatorname{Ass}(\mathbb{M})}\dim(\mathrm{A}/\mathfrak{p})=\dim(\mathrm{M})\) (VIII, § 1, 第 4 节, 注 2)，并且重新得到不等式 \(\operatorname{prof}(\mathrm{M})\leqslant\dim(\mathrm{M})\) （对 \(\mathrm{M}\neq0\) ）（第 4 节, 定理 2 的推论 2）。

